\section{Conclusion}

We connect a known isotropic ZGV critical ring to its known anisotropic
isolated-point endpoint through a coefficient-resolved geometric and temporal
reduction.  Full-elastic sensitivities determine \(V_4\) and \(B\), while the
exact determinant, independent spectral discretization, and centered
finite-difference checks close the coefficient evidence.  The controlled
normal form gives the exact critical-point cardinality for sufficiently small
perturbations, while the finite-resolution calculation resolves the corresponding alternating
minima and saddles with the predicted splitting, radial shift, remainder, and
sign reversal.  Response quadrature that passes the inter-resolution
phase-discrepancy test then connects the Bessel transition to the exact-Hessian
fixed-anisotropy Morse endpoint.  This
connection explains how the change in stationary-set dimension and Hessian rank
adds angular stationary phase and changes the computed temporal power law.  The
response comparison is performed without fitted alignment of amplitude, phase,
frequency, or time.

The two temporal descriptions express noncommuting limits.  In the joint
limit \(\varepsilon\to0\), \(t\to\infty\), with
\(\tau=\varepsilon V_4t=O(1)\), the radially localized ring retains its
\(t^{-1/2}\) factor and acquires the \(J_0(\tau)\) modulation.  At fixed,
resolved nonzero anisotropy, the isolated Cartesian Morse points instead give
the ultimate \(t^{-1}\) signed-Hessian sum, whose remainder is not uniform as
the ring is restored.  The exact minimum--saddle critical-frequency separation
is twice the modulation angular-frequency magnitude; keeping those quantities
distinct prevents a frequency-feature separation from being confused with the
rate that multiplies time.

The validity scope is restricted to one simple tracked branch of a lossless
infinite plate, a declared local annulus, a nonnodal source and weight class,
and the controlled weak-anisotropy family.  The results are computation-only
and make no global count across all branches; the finite-anisotropy silicon
stress test lies outside the asymptotic evidence.
Weak loss, finite boundaries, layered or defective plates, and near-closing
eigengaps require extensions of the coefficient and error analysis.
Nonlinear amplitude equations are a separate problem and belong to a distinct
study rather than to the linear response developed here.
